% =========================================================================
% TRANSFORMAÇÕES AVANÇADAS E SÍNTESE DE CUSTO MÍNIMO
% =========================================================================
\chapter{Transformações Avançadas e Síntese de Custo Mínimo}\label{sec:mcf_comparacao}

\section{A Transformação Min-Cost Max-Flow}
\label{sec:mcmf}

Uma variante comum do MCF é o problema do \textbf{Fluxo Máximo de Custo Mínimo} (\textit{Min-Cost Max-Flow}): dentre todos os fluxos máximos de $s$ a $t$, encontrar aquele com o menor custo total. Este problema pode ser reduzido ao MCF padrão com demandas através de uma construção simples (veja~\cite{ahuja1993network}).

A transformação consiste em adicionar à rede original um \textbf{arco artificial} $(t, s)$ do sumidouro para a fonte com capacidade infinita ($c(t,s) = \infty$) e custo suficientemente negativo ($w(t,s) = -M$, onde $M > \sum_{(u,v) \in A} |w(u,v)| \cdot c(u,v)$). Em seguida, define-se $b(v) = 0$ para todos os vértices, transformando o problema em um MCF puro sem demandas.

O custo negativo no arco artificial incentiva o algoritmo a maximizar o fluxo circulante pelo ciclo $s \leadsto t \to s$, garantindo simultaneamente o fluxo máximo e o custo mínimo. Qualquer algoritmo de MCF (incluindo o \textit{Network Simplex}) aplicado a esta rede transformada produz a solução desejada.

\paragraph{O Network Simplex como Motor Unificado para Fluxo Máximo Puro.}
O problema do fluxo máximo pode ser formulado e resolvido diretamente pelo \textit{Network Simplex} através desta modelagem. Para tanto, basta definir o custo de todos os arcos da rede original como zero ($w(u,v) = 0, \forall (u,v) \in A$) e atribuir custo estritamente negativo $w(t,s) = -1$ ao arco de retorno $(t,s)$ com capacidade infinita ($c(t,s) = \infty$), com demandas nodais nulas ($b(v) = 0$).

Sob esta formulação, a função objetivo do Simplex em Redes:
\begin{equation}
    \min \sum_{(u,v) \in A} w(u,v) f(u,v) = \min \big( -1 \cdot f(t,s) \big) \equiv \max f(t,s)
\end{equation}
minimiza o custo negativo do arco artificial, o que equivale a maximizar o fluxo líquido de $s$ a $t$ que retorna por $(t,s)$.

A estrutura desta transformação em rede de circulação com arco artificial de retorno é ilustrada na Figura~\ref{fig:circulacao_retorno}.

\begin{figure}[!ht]
\centering
\begin{tikzpicture}[thick]
    \node[source node] (s) at (0, 1.5) {$s$};
    \node[flow node] (u) at (3.5, 3.0) {$u$};
    \node[flow node] (v) at (3.5, 0.0) {$v$};
    \node[sink node] (t) at (7.0, 1.5) {$t$};

    \node[font=\footnotesize, left=3pt] at (s.west) {$b(s)\!=\!0$};
    \node[font=\footnotesize, above=3pt] at (u.north) {$b(u)\!=\!0$};
    \node[font=\footnotesize, below=3pt] at (v.south) {$b(v)\!=\!0$};
    \node[font=\footnotesize, right=3pt] at (t.east) {$b(t)\!=\!0$};

    \draw[flow edge] (s) -- node[edge lbl above] {$c\!=\!10, \; w\!=\!0$} (u);
    \draw[flow edge] (s) -- node[edge lbl below] {$c\!=\!8, \; w\!=\!0$} (v);
    \draw[flow edge] (u) -- node[edge lbl, right=3pt] {$c\!=\!5, \; w\!=\!0$} (v);
    \draw[flow edge] (u) -- node[edge lbl above] {$c\!=\!7, \; w\!=\!0$} (t);
    \draw[flow edge] (v) -- node[edge lbl below] {$c\!=\!10, \; w\!=\!0$} (t);

    \draw[tree edge, red, line width=1.8pt, dashed, bend left=60, looseness=1.4] (t) to node[below=5pt, font=\footnotesize, text=red!80!black, fill=white, inner sep=2pt] {Arco de Retorno Artificial: $c(t,s)\!=\!\infty, \; w(t,s)\!=\!-1$} (s);
\end{tikzpicture}
\caption{Redução do Problema do Fluxo Máximo a Fluxo de Custo Mínimo através de Circulação. O arco artificial de retorno $(t,s)$ com custo $w(t,s) = -1$ e capacidade infinita induz a saturação da circulação pelo ciclo $s \leadsto t \to s$, maximizando $f(t,s) = |f|$ com demandas nodais nulas ($b(v)=0$).}
\label{fig:circulacao_retorno}
\end{figure}

\paragraph{Justificativa Teórica e de Projeto sobre o Max-Flow Network Simplex.}
A literatura clássica (Ahuja, Magnanti e Orlin~\cite{ahuja1993network}, Goldfarb e Hao~\cite{goldfarb1990primal} e Tarjan~\cite{tarjan1991efficiency}) descreve a especialização do algoritmo Simplex em redes para fluxo máximo, manipulando árvores geradoras com capacidades sem custos. No entanto, não implementamos uma classe dedicada \lstinline{MaxFlowNetworkSimplex} em \lstinline{FlowNetwork} pelos seguintes motivos:
\begin{enumerate}
    \item \textbf{Complexidade e Desempenho em Redes sem Custos:} Em redes sem custos, algoritmos combinatórios dedicados são assintoticamente superiores: o algoritmo de Dinic opera em tempo $\mathcal{O}(|V|^2 |A|)$ (e $\mathcal{O}(|V|^{1/2} |A|)$ em redes unitárias), e o Push-Relabel com \textit{Highest Label} e \textit{Gap Heuristic} atinge $\mathcal{O}(|V|^2 \sqrt{|A|})$. O Simplex em redes não possui garantia polinomial no pior caso e requer a manutenção dinâmica da árvore geradora básica a cada iteração, gerando sobrecarga computacional na ausência de custos.
    \item \textbf{Equivalência Algorítmica e Modularidade:} O \textit{Max-Flow Network Simplex} e o \textit{Network Simplex} com custos nulos e arco artificial realizam sequências de pivoteamento idênticas, com custo de $\mathcal{O}(|V|)$ por iteração e mesmo limitante de pior caso $\mathcal{O}(|V|^2 \cdot |f^*|)$. A especialização apenas simplificaria o cálculo de potenciais nodais para uma partição booleana, mantendo idêntica a complexidade assintótica. Sob a ótica de engenharia de software, criar uma classe derivada exigiria duplicar o gerenciamento da base em árvore da classe \lstinline{NetworkSimplex}, violando o princípio DRY sem ganho prático de desempenho.
\end{enumerate}
Assim, a redução via arco de retorno artificial resolve o problema utilizando a implementação existente de \lstinline{NetworkSimplex}.


\section{Síntese Comparativa dos Três Paradigmas de Custo Mínimo}

A Tabela~\ref{tab:min_cost} sintetiza os paradigmas, invariantes mantidas e as complexidades assintóticas dos três algoritmos de fluxo de custo mínimo desenvolvidos nesta Iniciação Científica.

\begin{table}[!ht]
\centering
\small
\setlength{\tabcolsep}{5pt}
\begin{tabular}{p{4.2cm} l p{7.5cm}}
\toprule
\multicolumn{1}{c}{\textbf{Algoritmo}} & \multicolumn{1}{c}{\textbf{Complexidade}} & \multicolumn{1}{c}{\textbf{Mecanismo / Estratégia}} \\
\midrule
\multicolumn{3}{c}{\bannertext{\hyperref[sec:cycle_canceling]{Abordagem Primal Direta: Cancelamento de Ciclos}}} \\
\midrule
\hyperref[ap:cycle_canceling]{Cycle Canceling} & $\mathcal{O}(|V| \cdot |A|^2 \cdot C \cdot W)$ & Mantém viabilidade primal e cancela ciclos residuais negativos via SPFA \\
\midrule
\multicolumn{3}{c}{\bannertext{\hyperref[sec:ssp]{Abordagem Dual Construtiva: Caminhos de Custo Mínimo}}} \\
\midrule
\hyperref[ap:successive_shortest]{SSP (SPFA)} & $\mathcal{O}(|V| \cdot |A| \cdot F)$ & Caminhos aumentantes mais curtos em custos residuais via Bellman-Ford/SPFA \\
\hyperref[ap:successive_shortest_dijkstra]{SSP (Dijkstra + $\pi$)} & $\mathcal{O}(F \cdot |A| \log |V|)$ & Potenciais nodais $\pi$ ($\bar{w}_\pi \ge 0$) permitindo Dijkstra com heap binária \\
\midrule
\multicolumn{3}{c}{\bannertext{\hyperref[sec:network_simplex]{Abordagem Primal-Dual em Árvore: Estrutura de Base}}} \\
\midrule
\hyperref[ap:network_simplex]{Network Simplex} & $\mathcal{O}(|V| \cdot |A| \cdot C \cdot W)^*$ & Pivoteamento local em árvores geradoras via ciclos fundamentais ($\mathcal{O}(|V|)$) \\
\bottomrule
\end{tabular}
\caption{Comparação dos algoritmos de Fluxo de Custo Mínimo ($C = \max c$, $W = \max |w|$, $F = |f^*|$). $^*$Embora o Network Simplex possua pior caso teórico exponencial, cada iteração custa $\mathcal{O}(|V|)$ e seu desempenho prático é superior ao dos métodos de aumentos sucessivos na maioria das instâncias práticas.}
\label{tab:min_cost}
\end{table}

A condição de inexistência de ciclos direcionados de custo negativo no digrafo residual $D_f$ rege as três abordagens sob diferentes estruturas:
\begin{itemize}
    \item \textbf{Cycle Canceling (Primal Direto):} Mantém a viabilidade primal do fluxo e cancela iterativamente ciclos negativos no digrafo residual;
    \item \textbf{Successive Shortest Path (Dual Construtivo):} Mantém a otimalidade de custos como invariante ao rotear fluxo exclusivamente por caminhos residuais de custo mínimo, atingindo a viabilidade por aumentos sucessivos de fluxo;
    \item \textbf{Network Simplex (Primal-Dual em Árvore):} Mantém soluções básicas viáveis sobre árvores geradoras. A inclusão de um arco não básico com custo reduzido negativo induz um ciclo fundamental de custo negativo, que é saturado durante o pivoteamento.
\end{itemize}

Enquanto o \textit{Cycle Canceling} e o \textit{Successive Shortest Path} realizam buscas diretas em digrafos residuais, o \textit{Network Simplex} explora a estrutura combinatória das bases em árvore. A correspondência entre bases do programa linear e árvores geradoras permite substituir operações matriciais por travessias lineares $\mathcal{O}(|V|)$ com os vetores \lstinline{parent}, \lstinline{depth} e \lstinline{tree_adj}, resultando em menor tempo de execução na prática.

\vspace{0.8cm}

